% ============================================
% DOC-01
% Gerado de molambudos.md — build_miolo.py
% ============================================
\section*{\textbf{DOC-01} --- A Exumação}
\markboth{DOC-01 --- A Exumação}{}

\noindent \textit{Notas do coveiro Genésio Barbosa}

\noindent \textbf{Notas do coveiro Genésio Barbosa}
\textit{Transcrição literal do depoimento prestado em 14 de junho de 1979 à direção do Hospital Colônia. Documento arquivado sob o código C-1260-79.}

\vspace{0.5em}\noindent\rule{\textwidth}{0.5pt}\vspace{0.5em}

\noindent Era a terceira cova do dia. A gente sabe quais são as antigas porque a terra assentou. Covas recentes, a terra fica fofa, dá pra cavar com as mãos quase. As antigas, não. A terra endurece, forma uma crosta por cima, e quando você enfia a pá parece que vai bater em pedra.

\noindent A do 1.260 era antiga. Muito antiga. Pelo mapa, 1917. Um dos primeiros enterrados ali.

\noindent Eu cavei. Cavei uns quarenta centímetros e a pá bateu em madeira. Cavei mais. O caixão --- era um caixão, mas não era. Sabe quando a madeira apodrece e vira terra? Não é que desaparece. Ela vira outra coisa. Uma pasta escura que ainda tem a forma do que foi.

\noindent Eu abri com a pá. Não por desrespeito. É que não dava pra puxar a tampa --- ela já era parte do chão.

\noindent O que eu vi dentro não era um esqueleto.

\noindent Era um corpo.

\noindent Não, não estou maluco. Era um corpo. A pele era escura, quase preta, mas estava \textit{inteira}. Como se tivesse sido enterrado na semana passada. O senhor pode perguntar pro Seu Nilo, que tava comigo. Ele viu também. A pele do rosto ainda cobria os ossos. A boca estava aberta. Os olhos ---

\noindent Os olhos não estavam lá. As órbitas estavam vazias. Mas dentro delas, no fundo, havia um brilho. Como se alguém tivesse acendido uma vela lá dentro e apagado. Não, não era vela. Era um reflexo. Amarelo. Muito fraco.

\noindent E no peito do corpo --- amarrado com uma tira de couro podre, apertada contra o osso --- tinha um caderno. O couro atravessava as costelas, por baixo do paletó, como se o homem tivesse sido enterrado com o livro colado ao peito. Eu precisei cortar a tira com a faca. O couro esfarelou na minha mão.

\noindent Eu sei que parece mentira. Um caderno enterrado com um paciente em 1917, amarrado ao corpo, sessenta e dois anos depois, com as páginas ainda \textit{úmidas}? Eu não acreditaria. Mas o senhor está segurando o caderno. O senhor viu as páginas. A tinta.

\noindent O senhor sabe que eu não estou mentindo.

\noindent Eu peguei o caderno. Quando toquei nele, senti um calor. Não era calor do sol. Era um calor úmido, como o de uma ferida coberta. Como se o caderno estivesse vivo ali dentro. Crescendo.

\noindent Abri a primeira página.

\noindent A letra era de criança. Tortinha, subindo e descendo, como quem aprendeu a escrever sozinho. Uma assinatura no topo:

\noindent \textit{Joaquim Antônio Correia --- Paciente 1.260}

\noindent E uma data.

\noindent \textit{1917 --- o ano que eu entrei aqui}

\noindent Mas o diário continuava depois. Página após página. O último registro era de 1979.

\noindent \textit{Três dias antes de eu desenterrar ele.}

\noindent O senhor entende? O homem estava morto há 62 anos. E o diário continuou sendo escrito. Na cova. No escuro. Com os dedos de um esqueleto.

\noindent Eu não dormi naquela noite.

\noindent Não dormi até hoje.

\noindent O senhor pode ficar com o diário. Eu não quero mais ele perto de mim.

\noindent --- \textit{Genésio Barbosa, 48 anos, coveiro do Hospital Colônia de Barbacena há 22 anos.}

\vspace{0.5em}
\noindent\textit{\textbf{↪ Rotas:} \rota{MEM-04} \rotasep \rota{LUC-02}}